\documentclass{article}
\usepackage{graphicx} % Required for inserting images
\usepackage{caption} % Required for \captionof
\usepackage{amsmath}
\usepackage{float}
\usepackage{comment}
\usepackage[T1]{fontenc}
\usepackage[utf8]{inputenc}
\usepackage{lmodern}
\usepackage{microtype}
\usepackage{xcolor}
\usepackage{framed}
\renewenvironment{leftbar}{%
  \def\FrameCommand{\textcolor{orange}{\vrule width 3pt} \hspace{10pt}}%
  \MakeFramed {\advance\hsize-\width \FrameRestore}}%
 {\endMakeFramed}

\title{Hőtani mérés}
\author{Kern Luca, Csongor Vincze}
\date{2026 Április 14}

\begin{document}

\maketitle

\section*{Absztrakt}

Mérési jegyzőkönyvünkben több, egymásra épülő hőtani mérés eredményeit ismertetjük. Kezdetben a digitális hőmérőinket kalibráltuk, majd ezután belekezdtünk az érdekfeszítőbb mérésekbe. A mérések pontossága érdekében meghatároztuk először a rendszerünk fontosabb paramétereit, mint például a rendszer hőveszteségét és a termosz hőkapacitását. Munkánkat nehezítette az, hogy a második mérési alkalommal másik termoszt kaptunk, mint az első alkalommal, így néhány mérést meg kellett ismételnünk. Ezen adatok ismeretében már el tudtunk kezdeni olyan dolgokat is vizsgálni, amik nem képezték a mérőrendszerünk szerves részét. Így többek között kimértük a víz, az alumínium és a réz fajhőjét is, végül pedig a jég olvadáshőjét is meghatároztuk.

\section*{1. Feladat: A hőmérsékleti adatok leolvasása digitális multiméterről}

Ebben a feladatban a multiméteres hőmérsékletleolvasással ismerkedtünk meg. A táblázatban szereplő együtthatók és a képlet alapján meghatároztuk, hogy a mért ellenállások mekkora hőmérsékletnek felelnének meg.

Alább találhatóak az adataink, amelyek azt mutatják be, hogy mekkorának mértük a szobahőmérsékletet a két különböző hőmérővel. Mint láthatjuk, ezek az adatok eltérnek. Ennek oka az, hogy az ellenállásos hőmérésnél egy eltolást okoz a kapcsolás soros ellenállása, ami hozzáadódik a hőmérő saját ellenállásához. Ez egy szisztematikus hiba. Ezt a következő feladatban tárgyaljuk részletesebben. Emellett természetesen a két Pt100 érzékelő sem teljesen azonos, és a kiértékeléshez használt képlet együtthatói is kismértékben eltérhetnek egymástól.
A mérést és kiértékelést a következőképpen folytattuk: Bedugtuk a multiméterbe a Pt100 szenzort, és az Easy DMM szoftverrel elindítottunk egy mérési sorozatot. Ezt kiexportáltuk egy csv fájlba. A csv fájlt beolvastuk, és a feladatban megadott képlet szerint (feketére: $ R[\Omega]=100 \Omega+0.385 \frac{\Omega}{{ }^{\circ} \mathrm{C}} \cdot T\left[{ }^{\circ} \mathrm{C}\right] $, zöldre (3-as termosz): $ R[\Omega]=100 \Omega+0.425 \frac{\Omega}{{ }^{\circ} \mathrm{C}} \cdot T\left[{ }^{\circ} \mathrm{C}\right] $) visszaszámoltuk a hőmérsékleteket, és ezt ábrázoltuk. Ezt mutatja az alábbi grafikon.

% zöld termosz:\\
% $R_{zold}=110.4 \Omega$\\
% $T_{zold}= \text{ 24.5°C}$\\
% fekete termosz:\\
% $R_{fekete}=111.3 \Omega$\\
% $T_{fekete}= \text{ 26.6°C}$\\


\begin{figure}[H]
    \centering
    \includegraphics[width=0.8\textwidth]{1_1_feladat.pdf}
    \caption*{1. feladat: Mért hőmérséklet az idő függvényében}
    \label{}
\end{figure}

A két platinahőmérő közel ugyanazt a hőmérsékletet mérte (27.8 és 27.9 °C), ami magasabb, mint a szobahőmérséklet, viszont ebbe még nincsen beleszámítva a hozzáadódó soros ellenállás. Amennyiben ezzel az értékkel kompenzálunk, a hőmérsékleti értékek csökkenni fognak.

\begin{leftbar}
\textbf{Megjegyzés:} \textit{Ebben a feladatban a labor alatt összekevertük a termoszok adatait, így a naplóban még más értékeket mutat a grafikon. Emellett itt levágtuk azt a részt, ahol kontakthiba (vagy más mérési zavaró tényező) miatt hibás értékeket mutatott a multiméter. Így csak az első 25 másodpercet ábrázoltuk, de ez bőven elég volt a hőmérők működésének validálására. Ahogy az ábra is mutatja, a fekete termosz ingadozása nagyon csekely, és a zöld termosz is megbízhatóan mért (az ábrázolt 25 s-os szakaszban kicsit emelkedett a zöld termosz szenzora által mért hőmérséklet, de ez igen kis mértékű volt. Utána valamilyen zavaró tényezőnek köszönhetően kicsit lecsökkent a zöld termosz hőmérséklete).}
\end{leftbar}

\section*{2. Feladat: A platinahőmérő kalibrálása}

Ebben a feladatban a hőmérők kalibrálását végeztük el, amire azért volt szükség, mert ahogy a hőmérők bekötésre kerültek, a saját ellenállásukhoz hozzáadódott egy soros ellenállás is a vezetékeknek és a banándugóknak köszönhetően. Ez a mért hőmérsékleti érték eltolódásához vezetett volna kalibrálás hiányában. A kalibrálást a következőképpen végeztük: mindkét termoszba hideg vizet töltöttünk úgy, hogy a víz ellepje a hőmérőt. Ezután a hőmérsékleti egyensúly beállta után megmértük a termoszokban a víz hőmérsékletét higanyos hőmérővel, és ugyanekkor megmértük a platinahőmérővel is. A mért hőmérsékletek különbségéből visszaszámoltuk a soros ellenállásokat.

\begin{leftbar}
\textbf{Megjegyzés:} \textit{A második alkalommal mi kaptuk meg a zöld termosz helyett a pirosat, így a második mérési alkalommal a piros termoszban lévő hőmérőhöz adódó ellenállást is meghatároztuk, ezért van három sor a táblázatunkban.}
\end{leftbar}

\begin{table}[H]
    \centering
    \renewcommand{\arraystretch}{1.2}
    \begin{tabular}{|l|c|c|c|}
        \hline
        Termosz & Hőmérővel mért hőmérséklet [°C] & Mért ellenállás [$\Omega$] & Soros ellenállás [$\Omega$] \\ \hline
        Zöld & 22,1 & 110,7 & 1,3 \\ \hline
        Fekete & 22,0 & 111,4 & 2,9 \\ \hline
        Piros & 22,4 & 110,29 & 1,7 \\ \hline
    \end{tabular}
    \caption{Termoszok hőmérsékleti korrekciója}
    \label{tab:hotan_2}
\end{table}

\begin{comment}
Termosz színkódja: zöld (3)\\
Egyensúlyi hőmérséklet: 22,1 °C\\
Mért ellenállás: 110,7 $\Omega$\\
Soros ellenállás: 1,3 $\Omega$\\

Termosz színkódja: fekete\\
Egyensúlyi hőmérséklet: 22,0 °C\\
Mért ellenállás: 111,4 $\Omega$\\
Soros ellenállás: 2,9 $\Omega$\\

Termosz színkódja: piros\\
Egyensúlyi hőmérséklet: 22,4 °C\\
Mért ellenállás: 110,29 $\Omega$\\
Soros ellenállás: 1,7 $\Omega$\\
\end{comment}

Az eredményeinkből látszik, hogy az egyes termoszokban lévő hőmérők valóban kicsit eltérően mérnek. Emellett az is megfigyelhető, hogy a soros ellenállások egymástól csak kismértékben térnek el.


\section*{3. Feladat: A termosz hőveszteségének mérése}

Ebben a feladatban a termosz hőveszteségét mértük. Megallapitottuk, hogy ha a zöld (3-as) termoszt megtöltünk forró vízzel, akkor mekkora a hőveszteség. Ennek a hőveszteségnek két fő oka lehet. Az egyik az, hogy maga a termosz is átmelegszik (ezt a következő feladatban tárgyaljuk részletesebben), illetve hogy a környezet felé is lead a rendszer valamennyi hőt. Mivel a víz és a termosz közti termikus egyensúly beállta után kezdtünk el mérni, így ezzel a méréssel azt tudtuk megállapítani, hogy a rendszer mennyi hőt ad át a környezetének. Körülbelül 5 percig gyűjtöttünk hőmérsékletadatokat, majd ábrázoltuk ezeket az adatokat az eltelt idő függvényében. Erre a grafikonra egyenest illesztettünk, és meghatároztuk a meredekségét, amiből kiszámítottuk, hogy egy perc alatt nagyjából mekkora a hőveszteség. Az egész mérés ideje alatt 0,91 °C-ot csökkent a hőmérséklet. Ebből visszaszámolva percenként (átlagosan) 0,14 °C-ot csökkent a hőmérséklet. A mért értékeket és az illesztett egyenest az alábbi grafikon mutatja.

\begin{figure}[H]
    \centering
    \includegraphics[width=0.8\textwidth]{3_1_graf.pdf}
    \caption*{3. feladat: A termosz hővesztesége az idő függvényében}
    \label{fig:3_1}
\end{figure}

A grafikonról leolvasható, hogy kicsi a hőveszteség, a meredekség szinte 0, ami megfelel a várakozásainknak.

\begin{leftbar}
\textbf{Megjegyzés:} \textit{Néhol láthatóak kisebb „lyukak” a mérésben. Ez az EasyDMM mintavételezése (sampling) miatt lehet. Már a nyers csv fájlban is megtalálhatóak néhol nagyobb időbeli ugrások.}
\end{leftbar}

% ! Ellenőrizze a későbbi mérések során a jellemzőtermalizációs időket, és ennek megfelelően vegye figyelembe a hőveszteségen alapulóhibát. Ellenőrizze, hogy ez a hiba jelentős vagy elhanyagolható.

\section*{4. Feladat: A termosz hőkapacitásának becslése}

%hibabecslés: nem a tömegmérés relatív hibája, hanem a 2-vel való osztásnak. Annak van egy 10-15-20%-os hibája. Ezt majd át kell gondolni!

Ebben a feladatban a vizsgált termoszok belső hőkapacitását becsültük meg geometriai alapon. Ahogy azt korábban már említettük, az a hőmennyiség, amit a mérőedény fala a mérés során felvesz, alapvetően befolyásolja az összes vizsgált folyamat energiamérlegét. Ehhez megmértük az egyes termoszok teljes tömegét. Mivel a felépítésük kétrétegű, feltételeztük, hogy a belső és a külső acélfaluk ugyanolyan vastag, és mivel közöttük kiváló hőszigetelő tulajdonságú vákuum található, nincsen számottevő hőcsere a két falrész között. Emiatt a termikus folyamatokban ténylegesen részt vevő hőtároló tömeget a termosz teljes tömegének felével közelítettük. Ezt követően az edények hőkapacitását a $C = c_{\text{acél}} \cdot \frac{m_{\text{termosz}}}{2}$ összefüggés alapján számítottuk ki, a rozsdamentes acél irodalmi fajhőjének ($c_{\text{rozsdamentes acél}} = 500 \text{ J/(kg}^{\circ}\text{C)}$) felhasználásával.

\begin{table}[H]
    \centering
    \begin{tabular}{|l|c|c|}
        \hline
        Termosz & Teljes mért tömeg [g] & Becsült hőkapacitás [J/$^{\circ}$C] \\\hline
        Zöld & 430 & 107,5 \\ \hline
        Fekete & 444 & 111 \\ \hline
        Piros & 380 & 95 \\ \hline
    \end{tabular}
    \caption{Termoszok hőkapacitásának becslése}
    \label{tab:hotan_4}
\end{table}

\begin{comment}
$m_{\text{fekete}}=444$ g\\
$m_{\text{zöld}}=430$ g\\
$m_{\text{piros}}=380$ g\\
$c_{\text{rozsdamentes acél}}=500 \text{ J/(kg}^{\circ}\text{C)}$\\
Becsült hőkapacitás:\\
$C_{\text{fekete}}=111$ J/°C\\
$C_{\text{zöld}}=107,5$ J/°C\\
$C_{\text{piros}}=95$ J/°C\\
\end{comment}

\section*{5. Feladat: A víz fajhőjének mérése elektromos fűtéssel}

%fontos: előtte és utána is mérni 1-2 percig!!!
Ebben a feladatban a víz fajhőjét határoztuk meg elektromos fűtéssel. A zöld termoszunkat megtöltöttük 0,77 liter hideg vízzel, majd a folyadék folyamatos belső hőleadása (elektromos fűtés) melletti melegedését vizsgáltuk. A mérés kezdetén feljegyeztük a termoszban lévő még fűtetlen víz tömegét és kezdőhőmérsékletét. Ahhoz, hogy a fűtőtestet a méréshez javasolt stabil 50 W-os hőteljesítményen tudjuk üzemeltetni, ismernünk kellett az ellenállását. A digitális multimétert ellenállásmérő módba kapcsolva az fűtőszál ellenállása $1,23\ \Omega$-nak adódott. A minél nagyobb biztonság, valamint a garantáltan stabil teljesítményleadás érdekében a fűtőtestet egy áramgenerátorra kötöttük, és finoman beállítottuk rajta a szükséges feszültséget. A leolvasott áramerősség végül 6,46 A, az ehhez tartozó feszültség pedig 7,95 V lett (melyből pontosan számítható a villamos teljesítmény). Ezen beállítások mellett pontosan 10 percig folyamatosan fűtöttük a vizet, közben mágneses keverőt alkalmazva a homogén hőmérséklet-eloszlás elérésére.

A rögzített hőmérsékletadatokat ábrázoltuk az eltelt idő függvényében. Az így kapott monoton növő szakaszra egyenest illesztettünk (ez volt az ahol bekapcsoltuk a fűtést), amiből meghatároztuk a hőfoknövekedés meredekségét. Végül a termodinamikai energiamérleg felírásával megállapítottuk a víz fajhőjét, ehhez a termosz korábban becsült hőkapacitását is felhasználtuk. A számolás eredményeként a víz fajhőjére $4533 \text{ J/(kg}^{\circ}\text{C)}$ adódott. Ami megint nem tökéletes, de nincs olyan távol az irodalmi $4178 \text{ J/(kg}^{\circ}\text{C)}$ értéktől. A relatív hiba 8,5 $\%$, ami egy ilyen mérési beállításnál elfogadható. 

\begin{comment}
$V_{\text{víz}}=$0,77 L\\
$m_{\text{víz}}=$774 g\\
$R_{\text{fűtőtest}}=1,23 \Omega$\\
$U_{\text{fűtőtest}}=$7,95 V\\
$I_{\text{fűtőtest}}=6,46$ A\\
$T_{kezdeti}=26,4$ °C\\
$T_{veg}=36,4$ °C\\
termosz hőkapacitása az előző feladatból\\
$c_{\text{víz}}=4092 \text{ J/(kg}^{\circ}\text{C)}$\\

Ezek alapján a fajhő $c_{\text{víz}}=4532,52 \text{ J/(kg}^{\circ}\text{C)}$
-> Mi ez az elentmondás?????
\end{comment}

\begin{figure}[H]
    \centering
    \includegraphics[width=0.8\textwidth]{4_1_graf.pdf}
    \caption*{5. feladat: Fűtött víz hőmérsékletfüggése}
    \label{fig:5_1}
\end{figure}



\section*{6. Feladat: Termosz hőkapacitásának mérése keveréssel}

Ebben a feladatban a melegvizes termosz hőkapacitását határoztuk meg kalorimetriás keveredési módszerrel (különböző hőmérsékletű víztömegeket megkeverve). A megfelelő termoszokba beleöntöttük a meleg, illetve a hideg vizet, majd precízen megmértük a tömegüket. A hidegvizes edényben 464 g, a melegvizesben pedig 467 g víz volt. Miután mindkét termoszban beállt a belső termikus egyensúly, leolvastuk a kezdeti hőmérsékleteket. Ezután -- a mágneses keverő folyamatos működtetése mellett -- óvatosan beleöntöttük a hideg vizet a melegvizes termoszba, mialatt a keverék hőmérsékletének változását is folyamatosan rögzítettük. Az így kapott mérési adatsorból grafikonokat készítettünk, meghatározva a termikus egyensúlyi értékeket. A termosz hőkapacitásának meghatározásához az alábbi összefüggést írtuk fel: \[
m_{\text{meleg víz}} \cdot c_{\text{v}} \cdot (T'_{\text{meleg víz}} - T'_{\text{közös}}) + C_{\text{termosz}} \cdot (T'_{\text{meleg víz}} - T'_{\text{közös}}) = m_{\text{hideg víz}} \cdot c_{\text{v}} \cdot (T'_{\text{közös}} - T'_{\text{hideg víz}}).
\]
A számítások elvégzéséhez feltételeztük, hogy a keverési folyamat ideje alatt a rendszer hőszigetelt, teljesen zártnak tekinthető (vagyis a környezetbe hőáramlások elhanyagolhatók), miközben a víz fajhőjének (szobahőmérsékleti) 4178 J/(kg$^{\circ}$C) irodalmi értékét vettük alapul. Ezen feltételezésekkel a zöld termosz vizsgált hőkapacitása $197,4 \text{ J/}^{\circ}\text{C}$-ra adódott. Ami ugyan több mint 2-szer akkora mint a 4. feladatban becsült érték, de a mérés pontatlanságához képest elfogadható. Itt nagyon nagy szerepet játszott az, hogy a mérést hosszú ideig végeztük, és így a rendszer nem hőszigetelt.

\begin{comment}
Ezek az első alkalom eredményei.\\
$m_{hideg}=464$ g\\
$m_{meleg}=467$ g\\
$T_{hideg}=110,4 \Omega$\\
$T_{meleg}=141,0 \Omega$\\
$T_{\text{egyensúlyi}}=125,0 \Omega$\\
$c_{\text{víz}}= 4178 \text{ J/(kg}^{\circ}\text{C)}$ \\  %irodalmi adat
$c_{\text{rozsdamentes acél}}=500 \text{ J/(kg}^{\circ}\text{C)}$\\
Zöld termosz hőkapacitása keveréssel: $C_{termosz}=197.4$\\
\end{comment}

Sajnos mivel a második laboratóriumi alkalommal a szokványos zöld helyett mi egy piros termoszt kaptunk (ami térfogatban is különbözött az eddigiektől), a kalorimetriás kalibrálást ismét meg kellett ejtenünk ezen a piros eszközön is.

A megismételt második alkalommal is teljesen ugyanezzel a metodikával dolgoztunk a piros termosznál, ekkor viszont 294 g hideg és 406 g meleg vizet öntöttünk egybe. A hideg víz kezdő hőmérséklete 22,3 $^{\circ}$C, míg a meleg vízé 85,0 $^{\circ}$C volt. A keveredés utáni egyensúly, 58,4 $^{\circ}$C értéken állt be. Viszont a fenti mérlegegyenlet visszaszámolásakor a piros termosz hőkapacitására csupán egy $23,6 \text{ J/}^{\circ}\text{C}$-os érték adódott, amely a valóságosnál vélhetőleg jóval kisebb, de legalább nem negatív. A mérést nem ismételtük meg harmadszor, a laborvezető instrukciója szerint.
\begin{comment}
Ezek a második alkalom eredményei.\\
$m_{hideg}=294$ g\\
$m_{meleg}=406$ g\\
$R_{hideg}=110,29 \Omega$\\
$R_{meleg}=134,45 \Omega$\\
$T_{hideg}=22,3$ °C\\
$T_{meleg}=85,0$ °C\\
$R_{\text{egyensúlyi}}=125,0 \Omega$\\
$T_{\text{egyensúlyi}}=58,4$ °C\\
$c_{\text{víz}}= 4178 \text{ J/(kg}^{\circ}\text{C)}$ \\  %irodalmi adat
$c_{\text{rozsdamentes acél}}=500 \text{ J/(kg}^{\circ}\text{C)}$\\
Piros termosz hőkapacitása keveréssel: $C_{termosz}=23,6$ J/K\\
\end{comment}
\begin{figure}[H]
    \centering
    \includegraphics[width=0.8\textwidth]{6_1_graf.pdf}
    \caption*{6. feladat: Zöld termosz hőkapacitásának mérése keveréssel}
    \label{fig:6_1}
\end{figure}

\begin{figure}[H]
    \centering
    \includegraphics[width=0.8\textwidth]{6_2_graf.pdf}
    \caption*{6. feladat: Piros termosz hőkapacitásának mérése keveréssel}
    \label{fig:6_1}
\end{figure}


\section*{7. Feladat: Az alumínium fajhőjének mérése}

Ebben a feladatban az alumínium fajhőjét határoztuk meg. Ehhez először is megmértük a kapott alumínium henger paramétereit, hogy térfogatának kiszámolásával megbecsülhessük a tömegét. A henger átmérője 40 mm-nek, a magassága pedig 150 mm-nek adódott. Ebből a számított térfogat 189 cm$^3$-nek adódott. Az alumínium sűrűségét 2,7 g/cm$^3$-nek véve a becsült tömeg 509 g lett. A tömeget lemértük mérleggel is, és így 534 g-ot kaptunk. Ez viszonylag nagyobb eltérés. A továbbiakban a mérleggel mért tömeget használtuk a számolásokban a pontosabb eredményekért. 

Ezek után mindkét termoszt feltöltöttük annyi vízzel (az egyiket hideggel, a másikat meleggel), hogy ha az alumínium hengert belelógatjuk, akkor teljesen elmerüljön benne, de a víz ne csorduljon ki (itt már a piros termoszt használtuk melegvizesként). A folyadékok tömegét megmértük, majd a hideg vízbe beletettük az alumínium hengert. Megvártuk, amíg beáll a termikus egyensúly, majd leolvastuk a kezdeti hőmérsékleteket. A kezdeti adatok a következőknek adódtak: a hideg víz tömege 415 g, hőmérséklete 25,45 $^{\circ}$C. A meleg termoszban lévő víz tömege szintén 415 g volt, kezdeti hőmérséklete 92,2 $^{\circ}$C. Ezek után a lehűlt hengert áttettük a melegvizes termoszba, és addig mértük a hőmérsékletadatokat, amíg az új termikus egyensúly be nem állt. Ezt az alábbi grafikonon ábrázoltuk. Az új egyensúlyi hőmérsékletet 75,8 $^{\circ}$C-nak mértük. A fent megadott adatokból és a végső egyensúlyi hőmérsékletből pedig kiszámoltuk az alumínium fajhőjét, a leírásban mutatott módszerrel. 
$m_{\text{v}} \cdot c_{\text{v}} \cdot (T'_{\text{v}} - T'_{\text{közös}}) + C_{\text{termosz}} \cdot (T'_{\text{v}} - T'_{\text{közös}}) = m_{\text{Al}} \cdot c_{\text{Al}} \cdot (T'_{\text{közös}} - T'_{\text{Al}}).$
Megint feltettük azt a számolások során, hogy a környezettel nincsen hőcsere, és megint a víz irodalmi fajhőjét használtuk (4178 J/(kg$^{\circ}$C)). A piros termosz hőkapacitását a közelített 95 J/K-nek vettük. 

% ! szamitas leirasa

% $d_{\text{henger}} = 40,0 \text{ mm}$\\
% $h_{\text{henger}} = 15 \text{ cm}$\\
% $V_{\text{henger}} = 188,5 \text{ cm}^3$\\
% $\rho_{\text{Al}} = 2,7 \text{ g/cm}^3$\\
% $m_{\text{becsült}} = 509 \text{ g}$\\
% $m_{\text{mért}} = 534 \text{ g}$\\
% A hideg termosz adatai:\\
% $m_{\text{víz}} = 415 \text{ g}$\\
% $T_{\text{hideg}} = 25,45^{\circ}\text{C}$\\
% A meleg termosz adatai:\\
% $m_{\text{víz}} = 415 \text{ g}$\\
% $T_{\text{víz}} = 92,2^{\circ}\text{C}$\\
% $T_{\text{egyensúlyi}} = 75,8^{\circ}\text{C}$\\
% $c_{\text{víz}} = 4178 \text{ J/(kg}^{\circ}\text{C)}$ \\  % irodalmi adat
% $C_{\text{termosz}} = 95 \text{ J/K}$\\
% $c_{\text{Al}} = 634 \text{ J/(kg}^{\circ}\text{C)}$\\

\begin{figure}[H]
    \centering
    \includegraphics[width=0.8\textwidth]{7_1_graf.pdf}
    \caption*{7. feladat: Mért hőmérséklet az idő függvényében}
    \label{fig:7_1}
\end{figure}

Az eredményünk ugyan nem lett tökéletes, de nem is volt annyira távol az irodalmi adattól. A mi mérésünk (és számolásunk) alapján az alumínium fajhője 634 J/(kg$^{\circ}$C)-nak adódott. Ez elég jelentős mérési hibával és bizonytalansággal terhelt, így 630 J/(kg$^{\circ}$C)-ra kerekítettünk. Az irodalmi adat 890 J/(kg$^{\circ}$C). Ez 29 $\%$-os hibát jelent. A mérés eredménye a körülményekhez képest elfogadható. A legnagyobb hibaforrás valószínűleg az volt, hogy a második alkalommal kapott piros termosz jóval kisebb volt, mint a feketék vagy a zöldek, így a hőmérő egészen az alumínium henger falához nyomódott, ez a homersekletmerest reszben befolyasolhatta. Emellett a piros termosz kis mérete miatt a keverés is nehézkesebben történt meg, így jóval pontatlanabb volt a hőmérséklet-mérés is. Nyilván valamennyi környezet felé történő hőszivárgás is fellépett, de ez nem volt számottevő: az egyensúlyi hőmérséklet beállta után a hőmérséklet már nem csökkent tovább, igazolva ezzel, hogy a mérés ideje alatt a hőveszteség kicsi volt (ez egybevág a 3. feladat eredményeivel is). A termosz hőkapacitásának becsült értéke viszont nagyon nagy hibával terhelt, így a relatív nagy eltérés egy jelentős része valószínűleg ebből a pontatlanságból adódott. A mérés hossza miatt itt is hibaforrásként szerepet játszik a hő környezetbe való távozása. Ez a harmadik feladat eredményei alapján a mérés hosszát tekintve akár a Celsius nagyságrandbe is eshet, ezzel jelentős hibát adva a mérésünk eredményének.
% !homersekleti korrekcio



\section*{8. Feladat: A sárgaréz fajhőjének mérése}

Ebben a feladatban az előzőhöz nagyon hasonló mérést végeztünk, azzal a különbséggel, hogy most nem alumínium, hanem sárgaréz fajhőjét állapítottuk meg. A réz henger paramétereit ugyanakkorának mértük, mint az alumínium hengerét. A továbbiakban is megegyezett a mérés menete az előző feladatéval. A mért és számított adatainkat az alábbiakban ismertetjük:

\begin{itemize}
    \item \textbf{A rézhenger paraméterei:}
    \begin{itemize}
        \item Átmérő: $d_{\text{henger}} = 40 \text{ mm}$
        \item Magasság: $h_{\text{henger}} = 15 \text{ cm}$
        \item Sűrűség: $\rho_{\text{henger}} = 8,73 \text{ g/cm}^3$
        \item Becsült tömeg: $m_{\text{becsült}} = 1646 \text{ g}$
        \item Mért tömeg: $m_{\text{mért}} = 1592 \text{ g}$
    \end{itemize}
    \item \textbf{A hideg termosz adatai:}
    \begin{itemize}
        \item Kezdeti hőmérséklet: $T_{\text{hideg}} = 26,2^{\circ}\text{C}$
    \end{itemize}
    \item \textbf{A meleg termosz adatai:}
    \begin{itemize}
        \item Víz tömege: $m_{\text{víz}} = 436 \text{ g}$
        \item Kezdeti hőmérséklet: $T_{\text{víz}} = 82,3^{\circ}\text{C}$
        \item Egyensúlyi hőmérséklet: $T_{\text{egyensúlyi}} = 66,5^{\circ}\text{C}$
    \end{itemize}
    \item \textbf{Fajhő és hőkapacitás adatok:}
    \begin{itemize}
        \item Víz fajhője (irodalmi): $c_{\text{víz}} = 4178 \text{ J/(kg}^{\circ}\text{C)}$
        \item Rozsdamentes acél fajhője: $c_{\text{rozsdamentes acél}} = 500 \text{ J/(kg}^{\circ}\text{C)}$
        \item Termosz hőkapacitása: $C_{\text{termosz}} = 95 \text{ J/K}$
        \item \textbf{Sárgaréz kiszámolt fajhője: $c_{\text{Cu}} = 472 \text{ J/(kg}^{\circ}\text{C)}$}
    \end{itemize}
\end{itemize}

\begin{figure}[H]
    \centering
    \includegraphics[width=0.8\textwidth]{8_1_graf.pdf}
    \caption*{8. feladat: Mért hőmérséklet az idő függvényében}
    \label{fig:8_1}
\end{figure}

A számítások elvégzéséhez a következő összefüggést használtuk: $m_{\text{v}} \cdot c_{\text{v}} \cdot (T'_{\text{v}} - T'_{\text{közös}}) + C_{\text{termosz}} \cdot (T'_{\text{v}} - T'_{\text{közös}}) = m_{\text{Cu}} \cdot c_{\text{Cu}} \cdot (T'_{\text{közös}} - T'_{\text{Cu}}).$ A sárgaréz fajhője 472 J/(kg$^{\circ}$C)-nak adódott, amelyet a bizonytalanságok tükrében kerekítve, 470 J/(kg$^{\circ}$C)-ként érdemes figyelembe venni. Ez a 385 J/(kg$^{\circ}$C)-os irodalmi adattal viszonylag jó egyezésben van, a fajhőt kb. 22\%-os relatív hibával határoztuk meg. Bár ez az eredmény ezúttal sem tökéletes, a körülményeket tekintve nem is nagyon távoli. A hibák lehetséges okai hasonlóak az előző mérésnél tapasztaltakhoz. A rézhengernek már fix felfüggesztése volt (nem kötél, hanem merev drót), így amikor belehelyeztük a melegvizes termoszba, még jobban akadályozta a mágneses keverő mozgását. Emellett a hőmérő továbbra is a henger oldalához szorult, így nem a folyadék átlagos hőmérsékletét mértük. Ezeken felül itt is nagyon jelentős hibaforrást jelent a termosz hőkapacitásának meglehetősen nagy bizonytalansága. A mérés hossza miatt itt is hibaforrásként szerepet játszik a hő környezetbe való távozása. Ez a harmadik feladat eredményei alapján a mérés hosszát tekintve akár a Celsius nagyságrandbe is eshet, ezzel jelentős hibát adva a mérésünk eredményének.

\begin{figure}[H]
    \centering
    \includegraphics[width=0.8\textwidth]{8_2_graf.pdf}
    \caption*{nema}
    \label{fig:8_2}
\end{figure}


\section*{9. Feladat: A jég olvadáshőjének mérése keveréssel}

Ebben a feladatban a jég olvadáshőjét próbáltuk megállapítani olyan módon, hogy a piros termoszunkba meleg vizet töltöttünk, és megvártuk, míg ebben a termoszban beáll a termikus egyensúly. Ekkor körülbelül 0 $^{\circ}$C-os hőmérsékletű jeget tettünk a meleg vízbe, és elkezdtük gyűjteni az adatokat. Ezeket később grafikon formájában ábrázoltuk. A mérést addig folytattuk, míg a termoszban lévő összes jég el nem olvadt. A mérés végén a termoszban lévő összes víz tömegének megmérésével meghatároztuk a jég tömegét. Ezekből az adatokból pedig kiszámoltuk a jég olvadáshőjét a következő összefüggés alapján: $L_{\text{j}} m_{\text{j}} + c_{\text{v}} \cdot m_{\text{j}}(T_{\text{közös}} - T_{\text{j}}) = c_{\text{v}} \cdot m_{\text{v}} \cdot (T_{\text{v}} - T_{\text{közös}}) + C_{\text{termosz}}(T_{\text{v}} - T_{\text{közös}}).$
 A számolásokhoz a víz fajhőjére az irodalmi adatot használtuk, és feltettük, hogy a rendszer és a környezet közti hőcsere elhanyagolható. A mért és számolt eredményeink alább láthatóak.

\begin{itemize}
    \item \textbf{Mért tömegek:}
    \begin{itemize}
        \item Víz tömege: $m_{\text{víz}} = 477 \text{ g}$
        \item Teljes tömeg (jég elolvadása után): $m_{\text{teljes}} = 510 \text{ g}$
        \item Jég tömege: $m_{\text{jég}} = 33 \text{ g}$
    \end{itemize}
    \item \textbf{Hőmérsékleti adatok:}
    \begin{itemize}
        \item Kezdeti hőmérséklet: $T_{\text{kezdeti}} = 93,2^{\circ}\text{C}$
        \item Egyensúlyi hőmérséklet: $T_{\text{egyensúlyi}} = 80,8^{\circ}\text{C}$
    \end{itemize}
    \item \textbf{Fajhő és állandók:}
    \begin{itemize}
        \item Víz fajhője (irodalmi): $c_{\text{víz}} = 4178 \text{ J/(kg}^{\circ}\text{C)}$
        \item Termosz hőkapacitása: $C_{\text{termosz}} = 95 \text{ J/K}$
        \item \textbf{A jég kiszámolt olvadáshője: $L_{f} = 450 \text{ kJ/kg}$}
    \end{itemize}
\end{itemize}

\begin{figure}[H]
    \centering
    \includegraphics[width=0.8\textwidth]{9_1_graf.pdf}
    \caption*{9. feladat: Jég olvadáshő}
    \label{fig:9_1}
\end{figure}

A hiba leginkább abból származhatott, hogy a jég hőmérsékletét nem tudtuk pontosan, és nem kizárt, hogy hidegebb volt, mint 0 $^{\circ}$C, így már az olvadás előtt is hőt vont el a víztől. Emellett természetesen megint szerepet játszott a termosz hőkapacitásának pontatlansága. És valamennyire a disszipálódhatott is hő, bár ez inkább a tized Celsius nagyságrandbe esik (a 3. feladat eredményei alapján).

\section*{Összegzés, Hivatkozás}

A méréseket elvégeztük, viszonylag tőrhető hibán belül. Eredményeinket ismertettük, ábrázoltuk és megadtuk a hibák lehetséges forrásait. Ahhogy a grafikonokból kiderül a grafikus hőmérsékleti korrekció nem javított volna lényegeset az eredményeinken mivel az elő és utóperiódus meredeksége a mérések során igen alacsony volt. Emellett a mért értékeink sehol nem lettek teljesen irreálisak. Így a mérésvezető enedélyével ezeket nem végeztük el.

Az alumínium és réz fajhőadatait innét vettük: https://gchem.cm.utexas.edu/data/section2.php?target=heat-capacities.php

víz fajhője: https://hu.wikipedia.org/wiki/V%C3%ADz_(adatlap)

\end{document}